%%%%%%%%%%%%%%%%%%%%%%%%%%%%%%%%%%%%%%%%%%%%%%%%%%%%%%%%%%%%%%%%%%%%%%
% Zero-Cost Supervision for Multilingual Text Anomaly Detection
% via Large Language Models
%
% Template: SBC (same as STIL 2025 paper)
% Status: SKELETON — see PAPER_ROADMAP.md for writing plan
%%%%%%%%%%%%%%%%%%%%%%%%%%%%%%%%%%%%%%%%%%%%%%%%%%%%%%%%%%%%%%%%%%%%%%

\documentclass[12pt]{article}

\usepackage{sbc-template}
\usepackage{graphicx,url}
\usepackage[utf8]{inputenc}
\usepackage{amsmath}
\usepackage{tabularx}
\usepackage{amssymb}
\usepackage{multirow}
\usepackage{booktabs}
\usepackage{siunitx}
\usepackage{algorithm}
\usepackage{algpseudocode}

\sloppy

\title{Zero-Cost Supervision for Multilingual Text Anomaly Detection\\via Large Language Models}

\author{Omitted for Double-Blind Review}

\address{Omitted for Double-Blind Review
}

\begin{document}

\maketitle

% ══════════════════════════════════════════════════════════════════════════════
\begin{abstract}
Anomaly detection can operate without labels, yet semi-supervised methods ---
which leverage even a handful of labeled instances --- consistently outperform
fully unsupervised ones. The bottleneck is annotation: anomalies are rare,
expensive to label, and often unavailable at deployment time.
We propose a zero-cost supervision pipeline that replaces human annotators with
a compact, open-weight LLM. Given an unlabeled corpus, the pipeline encodes
documents, samples a small candidate set, queries the LLM for weak anomaly
labels, and optionally fine-tunes the encoder contrastively before training a
shallow anomaly detector. On four datasets in English and Portuguese
(hate speech and topic classification), gains are strongest where labels are
hardest to obtain: on the two hate speech datasets, where unsupervised baselines
barely exceed chance (0.56--0.58 ROC-AUC), the pipeline recovers 54--67\% of
the gap to fully supervised models. Globally, mean ROC-AUC improves from 0.708
to 0.819, using 15--78$\times$ fewer confirmed anomaly labels.
\end{abstract}


% ══════════════════════════════════════════════════════════════════════════════
\section{Introduction}

Anomaly detection (AD) can be performed without any labels
in its fully unsupervised form, yet semi-supervised methods --- which exploit even
a handful of labeled anomaly instances --- consistently yield superior
performance~\cite{pang2021deep, xu-etal-2023-comparative}.
The catch is fundamental: anomalies are rare by definition, so labeled
examples are scarce in almost every real-world setting, regardless of how
well the anomaly concept is understood.
Text AD is inherently difficult because anomalies are rare, diverse, and often
subtly different from normal instances~\cite{Chandola2009AnomalyDetection, pang2021deep}.
Most existing benchmarks focus on English corpora and structured domains,
leaving multilingual and low-resource settings largely underexplored~\cite{Boutalbi2023, xu-etal-2023-comparative}.
A prior benchmark~\cite{stil} confirmed that access to ground-truth labels for as
few as 5\% of training documents consistently improves ROC-AUC across English
and Portuguese datasets --- but obtaining even that small annotation set is
precisely the bottleneck in practice.

The bottleneck is annotation --- and not only because it is expensive.
For newly deployed systems, emerging threat patterns, or low-resource language
communities, labeled anomaly examples are not just costly to obtain: they
literally do not exist at deployment time. Even active learning, which reduces
per-label cost by directing attention to the most informative
examples~\cite{settles2009active}, presupposes a human expert willing and
able to label on demand. Large language models (LLMs) have recently 
demonstrated the ability to match or approximate human-level annotation quality 
on a variety of NLP tasks~\cite{wang2021want, he2023annollm}, raising the question 
of whether they can replace human annotators entirely for anomaly detection.

We propose a zero-cost supervision pipeline for text AD that replaces the human 
annotator with an LLM. Given an unlabeled corpus, our pipeline: (i) encodes all 
documents with a pre-trained multilingual sentence encoder; (ii) selects a 
small candidate set via random or diversity-based sampling; (iii) queries the 
compact open-weight SLM Qwen2.5-7B-Instruct~\cite{qwen2025} (run locally in
4-bit quantization) with a structured,
task-specific prompt to assign a weak anomaly label to each candidate; and (iv) optionally 
fine-tunes the encoder using SetFit~\cite{tunstall2022efficient} before training a 
semi-supervised MLP detector on the full corpus. No human labels are used at any 
stage.

We evaluate across four datasets spanning two tasks and two languages:
two hate speech datasets (English and Portuguese) and two topic classification
datasets (English and Portuguese). We compare a one-class detector
(DeepSAD~\cite{Ruff2019}) and a discriminative classifier (MLP) under two
encoder configurations (with and without contrastive fine-tuning). 

Our best configuration yields the strongest gains on the two hate speech
datasets --- where unsupervised baselines barely exceed chance
(0.56--0.58 ROC-AUC) --- recovering 54--67\% of the performance gap to 
fully supervised models trained with ground-truth anomaly labels, and raising AUC to 0.73--0.83.
Globally, mean ROC-AUC improves from 0.708 (unsupervised) to 0.819. On 20~Newsgroups, the unsupervised baseline is already strong (0.920) and the gain is modest (22\%). 
On WikiNews, the best pipeline variant recovers 30\% of the gap to fully supervised, but SetFit brings no gain due to LLM annotation failures. The pipeline uses 15--78$\times$ fewer confirmed anomaly labels than a fully supervised model, 
which is trained with ground-truth labels for approximately 5\% of the training set. 

We further identify two failure modes: (i) task-level misalignment, where the LLM consistently fails to flag any anomaly, 
rendering contrastive fine-tuning inactive; and (ii) noisy fine-tuning, where contrastive learning 
on very few mislabeled positives degrades the embedding quality.


Our contributions are:
\begin{itemize}
    \item A fully automated, zero-cost annotation pipeline that 
          replaces human annotators with an LLM for semi-supervised text AD.
    \item Empirical evidence that a compact open-weight LLM (Qwen2.5-7B,
          deployable locally in 4-bit quantized form, $\sim$4\,GB VRAM)
          raises ROC-AUC by up to $+$0.254 on hate speech datasets where
          unsupervised methods fail, closing up to 67\% of the gap
          to fully supervised models.
    \item Cross-lingual evaluation on four datasets spanning two tasks
          (hate speech, topic classification) and two languages (English, Portuguese),
          demonstrating that a single pipeline generalises
          across language and task-difficulty axes.
    \item Identification of two failure modes that guide when LLM-based 
          supervision is and is not effective.
\end{itemize}

We structure the paper around four research questions:

\begin{itemize}
    \item \textbf{RQ1:} Can LLM-generated labels improve anomaly detection performance?
    \item \textbf{RQ2:} Do LLM-generated labels outperform purely unsupervised methods?
    \item \textbf{RQ3:} Does contrastive fine-tuning help when labels are noisy and scarce?
    \item \textbf{RQ4:} How label-efficient is the proposed pipeline compared to fully supervised models?
\end{itemize}

% ══════════════════════════════════════════════════════════════════════════════
\section{Related Work}

\subsection{Text Anomaly Detection}

Anomaly detection has been extensively studied for tabular and image data; 
for text, the field is less mature~\cite{pang2021deep, Boutalbi2023}. 
Early approaches apply one-class classifiers such as OC-SVM~\cite{OCSVM} 
or Isolation Forest~\cite{IForest} on bag-of-words or TF-IDF features. 
Neural approaches include CVDD~\cite{ruff-etal-2019-self}, which adapts 
one-class deep learning to text using self-attention over contextual 
word vectors, and DATE~\cite{manolache2021date}, which trains a detector 
via self-supervised token-level corruption. Xu et al.~\cite{xu-etal-2023-comparative} 
benchmark 22 algorithms across 17 corpora and find that pre-trained 
transformer encoders dramatically improve detection quality by providing 
richer representations. Das et al.~\cite{das2023few} extend deviation 
networks to few-shot text AD, showing that even a handful of labeled 
anomalies provides large gains. Ait-Saada \& Nadif~\cite{aitsaada2023} 
address the multi-topic short-text setting that is close to our datasets.

A recurring theme across this literature is that labeled anomalies are 
treated as given. Our work follows the few-shot semi-supervised paradigm 
but removes the human annotation assumption entirely, replacing it with 
LLM-generated weak labels.

\subsection{LLMs for Anomaly and OOD Detection}

Recent work explores whether LLMs can serve directly as anomaly detectors 
via zero-shot prompting. Liu et al.~\cite{liu2024llmood} evaluate several 
LLMs on out-of-distribution (OOD) detection tasks and show that LLMs 
achieve competitive performance on some benchmarks but fail to generalize 
across domains. Yang et al.~\cite{yang2025adllm} conduct a systematic 
benchmark (AD-LLM) of LLMs as detectors and find that their effectiveness 
is highly task-dependent. A broader survey by Xu \& Ding~\cite{xu2024survey} 
catalogs how LLMs have been applied across AD and OOD tasks, distinguishing 
between using LLMs as feature extractors, detectors, and label generators.

Our approach differs from all of the above: rather than using the LLM 
as the final detector at inference time, we use it solely as a (noisy) 
\textit{annotator} during training, then hand off to a classical 
semi-supervised AD model. This decoupling makes the inference cost identical 
to a standard MLP and avoids the brittleness of LLM-based decision boundaries.

\subsection{LLMs as Annotators}

Wang et al.~\cite{wang2021want} demonstrate that GPT-3 can annotate text 
for classification tasks at a fraction of the cost of crowdsourcing, while 
maintaining competitive label quality. He et al.~\cite{he2023annollm} propose 
AnnoLLM, a structured prompting framework that further improves annotation 
consistency. These works focus on standard classification, where the label space 
is well-defined. We extend the paradigm to anomaly detection, where 
\textit{what constitutes an anomaly} must be implicitly conveyed through 
the task-specific prompt, without explicit class boundaries.

\subsection{Few-Shot Contrastive Fine-Tuning}

SetFit~\cite{tunstall2022efficient} fine-tunes a sentence transformer 
using contrastive (Siamese network) learning from a small set of labeled 
pairs, without requiring a full classification head. It achieves strong 
few-shot performance with as few as 8 examples per class. We integrate 
SetFit as an optional embedding refinement step in our pipeline, 
conditioned on the LLM producing at least 5 putative anomaly examples. 
Our experiments reveal that the benefit of SetFit is conditional on label 
reliability and quantity, contributing a new finding to the understanding 
of SetFit under label noise.


% ══════════════════════════════════════════════════════════════════════════════
\section{Methodology}

\subsection{Problem Formulation}

\paragraph{Anomaly detection vs.\  out-of-distribution detection.}
The terms \textit{anomaly detection} (AD) and \textit{out-of-distribution (OOD) 
detection} are often used interchangeably in NLP, but they differ in emphasis.
OOD detection asks whether a test instance comes from a distribution different 
from the training distribution, and is typically studied in a one-class or 
open-set recognition setting where no anomaly examples are seen during training. 
AD is the broader paradigm: it encompasses unsupervised, semi-supervised, and 
supervised regimes and does not require the anomaly to be a fully distinct 
distribution. In the \textit{semi-supervised} AD variant we study here, a small 
set of labeled anomalies is available at training time to guide the detector, 
making the task strictly easier than pure OOD detection but still challenging 
due to class imbalance and the diversity of anomalous patterns.



\paragraph{Formal setting.}
Let $X_{\text{train}} = \{x_1, \dots, x_M\}$ be a training corpus drawn predominantly
from a normal class, with a small fraction $\rho \approx 5\%$ of anomalous instances
whose identities are unknown at training time. We assume no human-labeled anomalies
are available. Instead, we use an LLM to automatically assign weak binary labels
to a small candidate set $S \subset X_{\text{train}}$, $|S| = N \ll M$.

For each document $x_i \in S$, a structured prompt is constructed as:
\begin{equation}
    P_i = T(x_i,\, d_{\text{task}},\, c_{\text{normal}},\, c_{\text{anomaly}})
\end{equation}
where $T(\cdot)$ is the prompt-building function, $d_{\text{task}}$ is a one-sentence
dataset description, $c_{\text{normal}}$ describes what counts as a normal instance,
and $c_{\text{anomaly}}$ specifies the anomaly criterion (all dataset-specific;
see Appendix~\ref{app:prompts}). The LLM processes $P_i$ and returns a continuous
anomaly score $s_i \in [0,1]$ together with a one-sentence explanation:
\begin{equation}
    (s_i,\, r_i) = f_{\text{LLM}}(P_i)
\end{equation}
The score is binarized at threshold $\tau_{\text{ann}} = 0.5$ to obtain a weak label:
\begin{equation}
    \hat{y}_i = \mathbf{1}[s_i \geq \tau_{\text{ann}}]
\end{equation}
The resulting weak anomaly set
$\hat{A} = \{x_i \in S : \hat{y}_i = 1\}$
is used to train the semi-supervised detector.
Ground-truth labels are used \textit{only} for final evaluation on the held-out test set.

\subsection{LLM Annotation Pipeline}

Figure~\ref{fig:pipeline} illustrates the full pipeline, which proceeds in six steps:

\begin{enumerate}
    \item \textbf{Embedding:} All documents in $X$ are encoded into dense vectors 
          $Z \in \mathbb{R}^{N \times d}$ using the pre-trained multilingual sentence 
          encoder \textit{distiluse-base-multilingual-cased-v2}~\cite{reimers2019sentence}.
    \item \textbf{Sampling:} A candidate set $S \subset X$ of size $N$ is selected 
          from the training pool using one of two strategies (Section~\ref{sec:sampling}).
    \item \textbf{LLM Annotation:} Each document $x_i \in S$ is presented to an 
          LLM with a dataset-specific structured prompt. The prompt includes 
          (i) a description of the task context, (ii) criteria for normal vs.\ 
          anomalous instances, and (iii) a request for a continuous anomaly score 
          in $[0, 1]$. Scores $\geq 0.5$ are binarized to 1 (anomaly). 
          We use \textit{Qwen2.5-7B-Instruct}~\cite{qwen2025} in 4-bit GGUF
          quantization (Q4\_K\_M, $\approx$4\,GB VRAM), executed locally via
          llama.cpp~\cite{llamacpp} on the available GPU.
    \item \textbf{SetFit Fine-Tuning (optional):} If the LLM identifies at least 5 
          anomalous samples in $\hat{A}$, the sentence encoder is fine-tuned 
          contrastively using SetFit~\cite{tunstall2022efficient} on the weak labels. 
          If this threshold is not met (e.g., when the LLM fails to detect 
          anomalies in the candidate set), fine-tuning is skipped and the 
          original distiluse embeddings are used unchanged.
    \item \textbf{Re-encoding:} All training documents are re-encoded with the 
          (optionally fine-tuned) encoder.
    \item \textbf{AD Model Training:} A semi-supervised anomaly detector is trained 
          on the full training set, using $\hat{A}$ as labeled anomaly examples.
\end{enumerate}

% TODO: insert \begin{figure} pipeline diagram — figs/pipeline.png
% \begin{figure}[ht]
% \centering
% \includegraphics[width=\linewidth]{figs/pipeline.png}
% \caption{...}
% \label{fig:pipeline}
% \end{figure}

Algorithm~\ref{alg:pipeline} provides a complete pseudocode description.

\begin{algorithm}[ht]
\caption{Zero-Cost LLM Supervision Pipeline}
\label{alg:pipeline}
\begin{algorithmic}[1]
\Require Training corpus $X_{\text{train}}$; test corpus $X_{\text{test}}$; LLM $\mathcal{M}$;
         task prompt $p_{\text{task}}$; budget $N$; sampling mode $\sigma$;
         SetFit threshold $\tau = 5$
\Ensure  Anomaly scores for $X_{\text{test}}$
\State $Z \gets \textsc{Encode}(X_{\text{train}})$ \Comment{distiluse-base-multilingual-cased-v2}
\State $S \gets \sigma(X_{\text{train}},\, Z,\, N)$ \Comment{random or diversity (k-means)}
\For{each $x_i \in S$}
    \State $s_i \gets \mathcal{M}(x_i,\, p_{\text{task}}) \in [0,1]$ \Comment{LLM anomaly score}
    \State $\hat{y}_i \gets \mathbf{1}[s_i \geq 0.5]$
\EndFor
\State $\hat{A} \gets \{x_i \in S : \hat{y}_i = 1\}$ \Comment{weak anomaly set}
\If{$|\hat{A}| \geq \tau$}
    \State $\textsc{Encoder} \gets \textsc{SetFit}(\textsc{Encoder},\, \hat{A},\, S \setminus \hat{A})$
    \State $Z \gets \textsc{Encode}(X_{\text{train}})$ \Comment{re-encode with fine-tuned encoder}
\Else
    \State $Z$ unchanged \Comment{fewer than $\tau$ anomalies found; skip fine-tuning}
\EndIf
\State Train MLP classifier on $(X_{\text{train}},\, Z,\, \hat{A})$
\State \Return $\text{MLP}.\textsc{Score}(\textsc{Encode}(X_{\text{test}}))$
\end{algorithmic}
\end{algorithm}

\subsection{Sampling Strategies}
\label{sec:sampling}

We evaluate two candidate selection strategies for $S$, treating the choice
as an ablation to assess sensitivity to selection policy:

\begin{itemize}
    \item \textbf{Random:} Uniformly sample $N$ instances from the training pool.
    \item \textbf{Diversity:} Partition the embedding space into $N$ clusters via
          $k$-means++ and select the centroid-nearest instance from each cluster ---
          a cluster-based representative selection strategy from the active learning
          literature~\cite{settles2009active,sener2018active}. Covering the embedding
          space uniformly exposes the LLM to a broader range of semantic patterns,
          which may increase the chance of surfacing rare anomalous instances in a
          corpus dominated by normal text.
\end{itemize}

\subsection{Anomaly Detection Models}
\label{sec:models}

\paragraph{Model families for semi-supervised text AD.}
Semi-supervised anomaly detection methods can be broadly classified along two axes:
\begin{itemize}
    \item \textbf{One-class (generative-boundary) models} learn a compact 
          representation of the \emph{normal} class — typically by minimising 
          the volume of a hypersphere that encloses normal embeddings — and flag 
          anything outside that boundary as anomalous. 
          Representative methods include Deep SVDD~\cite{Ruff2018} and its 
          semi-supervised extension DeepSAD~\cite{Ruff2019}, which uses labeled 
          anomalies to explicitly push the boundary away from known anomalous 
          points. The one-class objective is well-motivated when anomaly 
          examples are very few and clean, because it does not require learning 
          the anomaly distribution.
    \item \textbf{Discriminative (binary-classification) models} directly model 
          $P(\text{anomaly}\,|\,x)$ from labeled positives and negatives. When 
          paired with high-quality pre-trained features, a shallow MLP trained 
          with binary cross-entropy can be competitive. Deviation 
          networks~\cite{Pang2019} (DevNet / DeepSAD variants) also belong to 
          this family, imposing a prior on the anomaly score distribution. 
          Discriminative models tolerate label noise better than one-class models 
          because they do not build a global boundary that can be distorted by a 
          single mislabeled point; instead, the cross-entropy loss averages out 
          the effect of a few noisy labels across many normal examples.
\end{itemize}
A third family — \textbf{self-supervised methods} such as DATE~\cite{manolache2021date} 
--- requires no labels at all but still benefits from fine-tuning; these are 
included in the unsupervised baselines of~\cite{stil} but are not a focus here 
because our goal is to measure the \emph{incremental} value of LLM-generated labels.

\paragraph{Model selection rationale.}
We choose one representative from each major paradigm: 
\textbf{DeepSAD} (one-class) and \textbf{MLP} (discriminative). 
DeepSAD is the standard benchmark for semi-supervised deep AD: it is the 
direct extension of Deep SVDD with labeled anomaly support and has been 
widely validated on structured and text data. The MLP is the simplest 
discriminative baseline and is known to be robust to moderate label noise. 
By comparing them under identical noisy-label conditions, we can answer 
whether the one-class inductive bias (\emph{assume anomalies are rare but diverse}) 
outperforms the discriminative inductive bias 
(\emph{learn what anomalies look like from examples}) when the examples themselves 
are noisy.

\paragraph{The $2 \times 2$ ablation design.}
Our proposed method, \textbf{MLP+SF}, uses a discriminative MLP classifier 
trained on SetFit-refined embeddings. To understand which components are 
responsible for the gains, we conduct a $2 \times 2$ ablation by 
independently disabling each design choice:
\begin{itemize}
    \item Disabling SetFit (using the raw pretrained embeddings) gives the 
          \textbf{MLP} variant, isolating the contribution of embedding 
          adaptation.
    \item Replacing the MLP with DeepSAD (keeping the raw embeddings) gives 
          the \textbf{DeepSAD} variant, isolating the contribution of model 
          architecture.
    \item Disabling SetFit \emph{and} using DeepSAD gives \textbf{DS+SF} 
          (DeepSAD with SetFit), completing the factorial design.
\end{itemize}
All four configurations share the same LLM-generated labels, so differences 
in results are attributable solely to the model and embedding choices, not 
to annotation randomness. Table~\ref{tab:2x2-design} summarises the design.

\renewcommand{\arraystretch}{1.0}
\begin{table}[ht]
\small
\centering
\caption{$2 \times 2$ ablation design. \textbf{MLP+SF} is our proposed method; 
         the remaining three cells are ablations that disable SetFit, 
         swap the classifier, or both.}
\label{tab:2x2-design}
\begin{tabular}{lcc}
\toprule
 & \textbf{Original embeddings} & \textbf{SetFit fine-tuned} \\
\midrule
\textbf{DeepSAD} (one-class)  & DeepSAD      & DS+SF \\
\textbf{MLP} (discriminative) & MLP          & \textbf{MLP+SF} (proposed) \\
\bottomrule
\end{tabular}
\end{table}

\subsection{Datasets and Evaluation}
\label{sec:datasets}

We evaluate on four datasets spanning two tasks and two languages,
building on the benchmark of~\cite{stil}. All datasets are adjusted to a 5\% anomaly
contamination rate in the training set. The majority class in each dataset is designated
as normal; all other classes form the anomaly set.

\renewcommand{\arraystretch}{0.9}
\begin{table}[ht]
\small
\centering
\caption{Dataset overview. All datasets are adjusted to 5\% anomaly contamination.}
\label{tab:datasets}
\begin{tabular}{lp{2.2cm}p{2.2cm}ccr}
\toprule
\textbf{Dataset} & \textbf{Normal class} & \textbf{Anomaly class} & \textbf{Task} & \textbf{Lang} & \textbf{Size} \\
\midrule
20 Newsgroups~\cite{lang1995newsweeder}          & sci.crypt posts   & comp.graphics posts   & TC & en & 1,048  \\
WikiNews~\cite{10.1371/journal.pone.0296929}     & Politics articles & Non-politics articles & TC & pt & 6,120  \\
Hate Speech Tweets~\cite{sharma2018hate}         & Non-hate speech   & Hate speech tweets    & HS & en & 31,205 \\
HateBR~\cite{DBLP:journals/corr/abs-2010-04543} & Non-offensive     & Offensive comments    & HS & pt & 3,675  \\
\bottomrule
\multicolumn{6}{l}{\scriptsize TC = Topic Classification; HS = Hate Speech detection.}
\end{tabular}
\end{table}

The four datasets form a 2$\times$2 grid of task and language.
They vary considerably in detection difficulty, as measured by the best unsupervised
baseline AUC (Table~\ref{tab:datasets}).
\textbf{20 Newsgroups} is the easiest: the anomaly class (comp.graphics threads)
is topically and lexically distant from the normal class (sci.crypt posts),
so a pretrained encoder already separates them ($\approx 0.92$ AUC unsupervised).
\textbf{HateBR} and \textbf{Hate Speech Tweets} are the hardest: hateful and
neutral text share the same social-media register, vocabulary, and style,
so discrimination requires knowledge beyond surface form ($\leq 0.58$ AUC unsupervised).
\textbf{WikiNews} sits in between (0.776): articles from different topic sections
partially overlap in vocabulary and journalistic style.
The harder datasets are the main target application for our pipeline,
since they are precisely where fully unsupervised baselines struggle most.

Performance is measured by ROC-AUC on the held-out test set. 
Results are averaged over 3 random seeds $\times$ 2 sampling 
strategies $\times$ 2 values of $N \in \{50, 200\}$, yielding 12 runs 
per configuration per dataset. We compare against two baselines: 
(i) the best unsupervised model from~\cite{stil} (trained with no labels), 
and (ii) fully supervised models (trained with ground-truth
labels for $\sim$5\% of the training set).


% ══════════════════════════════════════════════════════════════════════════════
\section{Experiments and Results}

\subsection{Experimental Setup}

All experiments were conducted on an NVIDIA Tesla T4 GPU (Google Colab). 
LLM annotation used Qwen2.5-7B-Instruct~\cite{qwen2025} in 4-bit GGUF
quantization (Q4\_K\_M), run locally via llama.cpp~\cite{llamacpp} on the same
T4 GPU at zero API cost. All four configurations (proposed \textbf{MLP+SF}
and the three ablation variants) share the same LLM-generated labels,
persisted to disk before the ablation runs, so that differences in results
reflect only model and embedding choices.
Each configuration was evaluated over 48 runs 
(4 datasets $\times$ 3 seeds $\times$ 2 sampling strategies $\times$ 2 values of $N$).

\subsection{RQ1: Do LLM Labels Improve Over Unsupervised Detection?}

Table~\ref{tab:unsup-baselines} contextualises the unsupervised ceiling across
datasets. Three representative methods plus the per-dataset best are shown;
full results for all seven methods appear in Appendix~\ref{app:unsup-detail}.
The key takeaway is a sharp difficulty divide: 20~Newsgroups is largely solved
without labels ($\geq 0.85$ for most methods), whereas the hate speech datasets
are barely above chance (HateBR: 0.561, Hate Speech Tweets: 0.575 --- best
method each). This is precisely where supervision is most needed and most
difficult to obtain.

\renewcommand{\arraystretch}{0.9}
\begin{table}[ht]
\small
\centering
\caption{ROC-AUC of representative unsupervised baselines from~\cite{stil}
         (full breakdown in Appendix~\ref{app:unsup-detail}).
         \textbf{Best}: highest AUC across all seven methods; used as the
         \textit{Unsup} reference in Table~\ref{tab:main-results}.}
\label{tab:unsup-baselines}
\begin{tabular}{lcccc}
\toprule
\textbf{Method} & \textbf{20 News} & \textbf{WikiNews} & \textbf{HS Tweets} & \textbf{HateBR} \\
\midrule
IForest      & 0.858 & 0.697 & 0.536 & 0.382 \\
AutoEncoder  & 0.902 & 0.757 & \textbf{0.575} & 0.474 \\
VAE          & \textbf{0.920} & \textbf{0.776} & 0.526 & 0.381 \\
\midrule
\textbf{Best} & \textbf{0.920} & \textbf{0.776} & \textbf{0.575} & \textbf{0.561} \\
\bottomrule
\end{tabular}
\end{table}

Table~\ref{tab:main-results} shows what LLM-generated weak supervision achieves on top of this ceiling.

\renewcommand{\arraystretch}{0.9}
\begin{table}[ht]
\small
\centering
\caption{Mean ROC-AUC per dataset. \textit{Unsup}: best of 7 unsupervised
         baselines from~\cite{stil} (full breakdown in Appendix~\ref{app:unsup-detail}).
         \textit{Sup.}: best fully supervised model with ground-truth labels~\cite{stil}.
         MLP, DeepSAD, DS+SF are ablation variants; \textbf{MLP+SF} is the proposed method.
         Gap$\uparrow$: fraction of the \mbox{[Unsup, Sup.]} interval recovered by the best pipeline variant (bold).}
\label{tab:main-results}
\begin{tabular}{lcccccc r}
\toprule
\textbf{Dataset} & \textbf{Unsup} & \textbf{DS+SF} & \textbf{DeepSAD} & \textbf{MLP} & \textbf{MLP+SF} & \textbf{Sup.} & \textbf{Gap$\uparrow$} \\
\midrule
20 Newsgroups      & 0.920 & 0.879 & 0.850 & \textbf{0.937} & \textbf{0.937} & 0.998 & 22\% \\
WikiNews           & 0.776 & 0.719 & 0.748 & \textbf{0.827} & 0.778 & 0.943 & 30\% \\
Hate Speech Tweets & 0.575 & 0.791 & 0.806 & 0.819 & \textbf{0.829} & 0.956 & 67\% \\
HateBR             & 0.561 & 0.620 & 0.595 & 0.672 & \textbf{0.731} & 0.873 & 54\% \\
\midrule
\textbf{Global}    & 0.708 & 0.752 & 0.750 & 0.814 & \textbf{0.819} & 0.942 & 43\% \\
\bottomrule
\end{tabular}
\end{table}

Our proposed \textbf{MLP+SF} surpasses the unsupervised ceiling on all datasets
and achieves the best pipeline AUC on three of four. The gains are most dramatic
on the two hate speech datasets: \textbf{HateBR} jumps from 0.561 to 0.731
(+0.170), and \textbf{Hate Speech Tweets} from 0.575 to 0.829 (+0.254) ---
recovering 54\% and 67\% of the performance gap between unsupervised and fully
supervised models respectively. These are precisely the datasets where practitioner
pain is highest: fully unsupervised methods are barely above chance,
yet the pipeline --- using only a local quantized LLM and no human annotation
--- almost halves the distance to fully supervised performance.

On \textbf{WikiNews} the \textbf{MLP} ablation (without SetFit) performs best
(0.827 vs.\ 0.778, Gap$\uparrow$ 30\%). Here the LLM struggles with the
task boundary --- politics vs.\ non-politics is a subtle distinction in short
news snippets --- producing few and noisy anomaly labels; contrastive
fine-tuning with those labels distorts the embedding space rather than
sharpening it. On \textbf{20 Newsgroups}, where the two topic classes are
lexically distant (cryptography vs.\ computer graphics), the unsupervised
ceiling is already high (0.920) and the pipeline adds a modest but consistent
gain (0.937, Gap$\uparrow$ 22\%) --- the LLM correctly flags almost no post
as anomalous, so SetFit is skipped entirely and \textbf{MLP+SF}
$\equiv$ \textbf{MLP}. Globally, the pipeline raises mean AUC from 0.708
to 0.819 with zero human annotation.

\subsection{RQ2: Which AD Model Is More Robust to Noisy Labels?}

Comparing the \textbf{MLP} and \textbf{DeepSAD} ablations --- which differ 
only in the classifier, holding embeddings fixed --- isolates the effect of 
architecture under noisy LLM labels. MLP consistently outperforms DeepSAD 
across all four datasets, with a global mean gain of $+$0.064. This is a 
counterintuitive result: one-class models are typically expected to be more 
robust when anomaly labels are noisy, since they primarily model the normal 
class. Instead, we find that the discriminative objective is more tolerant: 
a cross-entropy loss averaged over thousands of clean negative examples is 
less sensitive to a handful of mislabeled positives than a hypersphere whose 
boundary can be materially shifted by even a few anomalous training points.

\subsection{RQ3: Does SetFit Help Under Noisy Labels?}

Table~\ref{tab:setfit-effect} reports the paired AUC difference 
\textbf{MLP+SF} $-$ \textbf{MLP} (same labels, same AD model, only the encoder 
differs), and Table~\ref{tab:setfit-skip} shows how often SetFit actually ran.

\renewcommand{\arraystretch}{0.9}
\begin{table}[ht]
\small
\centering
\caption{SetFit effect: paired AUC difference MLP+SF $-$ MLP (same labels; encoder is the only variable).}
\label{tab:setfit-effect}
\begin{tabular}{lcc}
\toprule
\textbf{Dataset} & $\Delta$ \textbf{mean} & $\Delta$ \textbf{std} \\
\midrule
20 Newsgroups      & $+$0.000 & $\pm$0.000 \\
WikiNews           & $-$0.049 & $\pm$0.099 \\
Hate Speech Tweets & $+$0.010 & $\pm$0.038 \\
HateBR             & $+$0.059 & $\pm$0.082 \\
\midrule
\textbf{Global}    & $+$0.005 & $\pm$0.078 \\
\bottomrule
\end{tabular}
\end{table}

\renewcommand{\arraystretch}{0.9}
\begin{table}[ht]
\small
\centering
\caption{SetFit execution status and mean AUC by dataset and $N$.}
\label{tab:setfit-skip}
\begin{tabular}{llcccc}
\toprule
\textbf{Dataset} & $N$ & \textbf{Ran} & \textbf{Skipped} & \textbf{AUC (ran)} & \textbf{AUC (skipped)} \\
\midrule
20 Newsgroups      & 50  & 0 & 6 & ---   & 0.936 \\
20 Newsgroups      & 200 & 0 & 6 & ---   & 0.939 \\
WikiNews           & 50  & 4 & 2 & 0.727 & 0.822 \\
WikiNews           & 200 & 6 & 0 & 0.798 & ---   \\
Hate Speech Tweets & 50  & 3 & 3 & 0.848 & 0.725 \\
Hate Speech Tweets & 200 & 6 & 0 & 0.872 & ---   \\
HateBR             & 50  & 1 & 5 & 0.719 & 0.650 \\
HateBR             & 200 & 6 & 0 & 0.801 & ---   \\
\bottomrule
\end{tabular}
\end{table}

Two failure modes emerge. First, \textbf{task misalignment}: for 20~Newsgroups, 
the LLM conflates the anomaly class (comp.graphics) with the normal class 
(sci.crypt) — it rarely identifies any post as anomalous ($\sim$1 per run), 
so SetFit is always skipped and \textbf{MLP+SF} $\equiv$ \textbf{MLP}. 
Second, \textbf{noisy fine-tuning}: for WikiNews, contrastive fine-tuning 
with only $\sim$15 noisy anomaly labels per run distorts the embedding space 
($-$0.049 vs.\ MLP). When enough reliable positives exist (HateBR, $N=200$: 
$\sim$8 per run), SetFit delivers a substantial benefit ($+$0.150 over runs 
where it was skipped).

\subsection{RQ4: Label Efficiency --- Pipeline vs.\ Fully Supervised}

Table~\ref{tab:label-efficiency} compares the number of confirmed anomaly labels
used by the fully supervised baseline versus the pipeline.

\renewcommand{\arraystretch}{0.9}
\begin{table}[ht]
\small
\centering
\caption{Label efficiency. Fully supervised (Sup.) receives ground-truth labels for $\sim$5\%
         of training; pipeline receives LLM-assigned labels from $N=200$ samples,
         of which only a fraction are flagged as anomalous.}
\label{tab:label-efficiency}
\begin{tabular}{lrrcr}
\toprule
\textbf{Dataset} & \textbf{Sup.} & \textbf{Pipeline} & \textbf{Ratio} & \textbf{AUC gap} \\
 & \textbf{(GT anom.)} & \textbf{($N$=200, mean)} & & MLP+SF vs Sup. \\
\midrule
20 Newsgroups      & 39   & $\sim$1  & $\sim$49$\times$ & $-$0.061 \\
WikiNews           & 233  & $\sim$15 & $\sim$15$\times$ & $-$0.165 \\
Hate Speech Tweets & 1189 & $\sim$15 & $\sim$78$\times$ & $-$0.127 \\
HateBR             & 140  & $\sim$8  & $\sim$18$\times$ & $-$0.142 \\
\bottomrule
\end{tabular}
\end{table}

Despite using 15--78$\times$ fewer confirmed anomalies than fully supervised models, the
pipeline achieves AUC gaps of only 6--17 pp. The most compelling case is
Hate Speech Tweets: with $\sim$15 LLM-identified anomalies versus 1,189 ground-truth
labels (a 78$\times$ disadvantage), the pipeline still closes 67\% of the
gap to fully supervised models. This demonstrates that LLM annotations carry 
meaningful supervisory signal even when sparse and noisy, substantially 
reducing annotation cost while preserving most of the performance gain 
of weak supervision.


% ══════════════════════════════════════════════════════════════════════════════
\section{Conclusion}

We presented a zero-cost supervision pipeline for multilingual text anomaly 
detection that eliminates the need for human annotators by leveraging LLMs 
to generate noisy training labels. Our study yields four main findings.

First, LLM-generated labels provide consistent improvements over unsupervised 
detection across all four datasets, raising the mean ROC-AUC from 0.708 to 
0.819 without any human annotation --- a gain that covers 31--67\% of the
gap to fully supervised models for three of the four datasets.

Second, the choice of anomaly detection model matters more than expected: 
standard MLP classifiers outperform DeepSAD by $+$0.064 global mean AUC 
under noisy LLM labels, suggesting that discriminative objectives are more 
tolerant of annotation noise than one-class objectives.

Third, SetFit contrastive fine-tuning is beneficial \textit{conditional} on 
label quality and quantity: it provides sizable gains when the LLM reliably 
identifies anomalies ($+$0.059 for HateBR), but fails when the LLM annotation 
is too sparse (20~Newsgroups, always skipped) or when few-shot noisy fine-tuning 
distorts the embedding space (WikiNews, $-$0.049 at $N=50$).

Fourth, the pipeline is remarkably label-efficient: despite using 15--78$\times$
fewer confirmed anomaly labels than fully supervised models, the AUC gap is only 6--17 pp.
The most striking example is Hate Speech Tweets, where a 78$\times$ label
disadvantage still yields 67\% of the gain achieved by fully supervised models over unsupervised.

For future work, we plan to investigate adaptive prompt optimization, larger 
LLMs with calibrated confidence scores, and light human-in-the-loop mechanisms 
to handle the failure modes identified here.

\section*{Acknowledgments}
\small
[Omitted for double-blind review.]
\normalsize

\bibliographystyle{sbc}
\bibliography{refs}

% ══════════════════════════════════════════════════════════════════════════════
\appendix
\section{Unsupervised Baseline Detail}
\label{app:unsup-detail}

Table~\ref{tab:unsup-full} reports the ROC-AUC of seven unsupervised methods
from~\cite{stil} on all four datasets. The \textbf{Best} row is used as the
\textit{Unsup} reference in Table~\ref{tab:main-results}.

\renewcommand{\arraystretch}{0.9}
\begin{table}[ht]
\small
\centering
\caption{Unsupervised baseline ROC-AUC from~\cite{stil}.}
\label{tab:unsup-full}
\begin{tabular}{lcccc}
\toprule
\textbf{Method} & \textbf{20 News} & \textbf{WikiNews} & \textbf{HS Tweets} & \textbf{HateBR} \\
\midrule
IForest          & 0.858 & 0.697 & 0.536 & 0.382 \\
LOF              & 0.856 & 0.676 & 0.498 & 0.561 \\
OCSVM            & 0.770 & 0.666 & 0.553 & 0.372 \\
DeepSVDD         & 0.702 & 0.591 & 0.548 & 0.515 \\
AutoEncoder      & 0.902 & 0.757 & \textbf{0.575} & 0.474 \\
VAE              & \textbf{0.920} & \textbf{0.776} & 0.526 & 0.381 \\
HBOS             & 0.917 & 0.743 & 0.540 & 0.377 \\
\midrule
\textbf{Best}    & \textbf{0.920} & \textbf{0.776} & \textbf{0.575} & \textbf{0.561} \\
\bottomrule
\end{tabular}
\end{table}

\section{Annotation Prompts}
\label{app:prompts}

All four datasets share the same prompt template; only the three context
fields ($d_{\text{task}}$, $c_{\text{normal}}$, $c_{\text{anomaly}}$) vary
per dataset. The template is:

\begin{small}
\begin{verbatim}
You are evaluating a text sample for anomaly detection.

Dataset context: {description}

Normal samples are: {normal_description}

Anomaly criterion (use this definition strictly):
{anomaly_criterion}

Rate the following text on a scale from 0.0 (clearly normal)
to 1.0 (clearly anomalous according to the criterion above).
Do not use any other notion of "anomalous".

Text:
"""
{text}
"""

Respond ONLY with a JSON object, no explanation outside the JSON:
{"anomaly_score": <float 0.0-1.0>,
 "reason": "<one sentence explaining your score>"}
\end{verbatim}
\end{small}

Table~\ref{tab:prompts} lists the three context fields for each dataset.

\renewcommand{\arraystretch}{1.1}
\begin{table}[ht]
\small
\centering
\caption{Per-dataset prompt context fields. $d_{\text{task}}$: task description;
         $c_{\text{normal}}$: definition of a normal instance;
         $c_{\text{anomaly}}$: anomaly criterion.}
\label{tab:prompts}
\begin{tabularx}{\textwidth}{lX}
\toprule
\textbf{Field} & \textbf{Content} \\
\midrule
\multicolumn{2}{l}{\textbf{20 Newsgroups}} \\
$d_{\text{task}}$ & English newsgroup posts. Decide if a post belongs to the comp.graphics newsgroup (anomalous, 1.0) or any other topic (normal, 0.0). \\
$c_{\text{normal}}$ & Posts about any topic other than computer graphics: sports, religion, politics, science, medicine, automobiles, space, etc. \\
$c_{\text{anomaly}}$ & Posts about computer graphics: image formats (GIF, JPEG, PNG), 3D rendering, raytracing, image processing, graphics software, display hardware. \\
\midrule
\multicolumn{2}{l}{\textbf{HateBR}} \\
$d_{\text{task}}$ & Brazilian Portuguese Instagram comments about politicians. Decide if a comment is offensive (anomalous, 1.0) or not offensive (normal, 0.0). \\
$c_{\text{normal}}$ & Opinions, criticism, or disagreement without personal attacks; political criticism, factual statements, neutral opinions, sarcasm without insults. \\
$c_{\text{anomaly}}$ & Comments flagged by a content moderator: direct insults, hate speech by race/gender/religion/orientation, homophobia, sexism, racism, or promotion of political violence (score $\geq 0.5$). \\
\midrule
\multicolumn{2}{l}{\textbf{Hate Speech Tweets}} \\
$d_{\text{task}}$ & English tweets. Decide if a tweet contains hate speech targeting an identity group (anomalous, 1.0) or not (normal, 0.0). \\
$c_{\text{normal}}$ & Tweets that do not attack people based on group identity: opinions, personal conflicts, profanity with no identity-group target. \\
$c_{\text{anomaly}}$ & Tweets attacking or dehumanizing people by race, ethnicity, gender, orientation, religion, nationality, or disability --- including coded language, irony, and derogatory generalizations (score $\geq 0.5$). \\
\midrule
\multicolumn{2}{l}{\textbf{WikiNews}} \\
$d_{\text{task}}$ & Portuguese WikiNews articles. Decide if the article is from the Politics section (normal, 0.0) or from any other section (anomalous, 1.0). \\
$c_{\text{normal}}$ & Articles whose primary topic is political: elections, government decisions, diplomacy, military conflicts, legislation, political figures, geopolitical crises. \\
$c_{\text{anomaly}}$ & Articles whose primary topic is clearly non-political: pure health/medicine, sports results, culture/entertainment, science/technology, or environment --- with no significant political angle. \\
\bottomrule
\end{tabularx}
\end{table}

\end{document}
